\section{An inspectable execution trace}\label{sec:C15:trace}
The inspected implementation supplies a concrete read-write-maintain path. \path{HyperSkillProvider.provide_memory} handles a memory request, obtains or generates a plan, invokes the hypergraph retrieval, and formats the resulting context. \path{take_in_memory} accepts a completed trajectory and performs outcome-conditioned extraction. In the core, \texttt{HyperGraph.retrieve} implements the task and subtask routes, while \path{maintain_graph} performs pruning, structural merging, orphan cleanup, and persistence~\cite{repo_hyperskill2026}. These are identifiable source functions at a pinned revision; they have been inspected, not executed here.

One feedback boundary deserves explicit validation. The provider reads the trajectory metadata field \path{is_correct} with a default of false. If a caller omits that field for an indeterminate outcome, the extraction path treats the episode as failure rather than as unknown~\cite{repo_hyperskill2026}. A wrapper claiming evidence-based consolidation should require a typed outcome $y\in\{0,1,\bot\}$, update outcome-dependent utility only for $y\ne\bot$, and retain unresolved trajectories separately. This is a proposed integration safeguard; it is not a claim that the paper's labeled evaluation omitted its outcomes.

Consider a constructed example with two past trajectories. One retrieved edge contains subtasks for locating a source, verifying its date, and comparing a claim, together with skills for checking primary evidence and preserving uncertainty. A second edge contains source-location and comparison subtasks plus a skill warning against treating a product announcement as a paper. A new research query matches both through its decomposition, even if the original task descriptions use different domain vocabulary. Co-occurrence ranks the shared evidence-checking skill above a skill appearing in only one edge, subject to the configured candidate pool.

This is an algorithmic illustration, not a benchmark observation. It makes two failure conditions explicit. First, if decomposition misses date verification, a relevant edge can fail to enter $E_{\mathrm{sub}}$. Second, if a generic skill appears in many retrieved edges, frequency can dominate a rare warning. The task route can recover some missed edges, but no guarantee follows from taking the union. Testing this example with controlled distractor edges would isolate a representation and ranking effect.

In the inspected core, the retrieval budget is divided between task-hyperedge candidates and subtask candidates. The subtask search overfetches candidates before filtering for the correct node type. Skills and mistakes are then ranked separately and limited separately. Consequently, the final guidance count and token count are not determined by a single nominal $k$~\cite{repo_hyperskill2026}. A cost comparison should inspect the actual formatted context rather than treating the retrieval parameter as the complete evidence budget.

The implementation also preserves instance-specific subtasks differently from abstract skills. Subtask insertion uses content-aware exact identity, while reusable abstractions can undergo semantic deduplication and structural merging. This distinction reflects the need to retain episode-specific planning context. It prevents one broad subtask label from automatically collapsing every occurrence, but it does not eliminate extraction or compatibility errors in the skill nodes.

A source-level reproduction should log the generated plan, retrieved subtask identifiers, task-route hyperedges, fused edge set, ranked skill and mistake identifiers, utility updates, and maintenance changes. These observables connect an outcome to the actual mechanism. Reporting only a final success rate and graph size would leave retrieval, extraction, and maintenance contributions difficult to distinguish.

